\section{Background and Related Work}
\label{sec:background}

The images generated by FLS can look deceptively similar to an optical camera, but the information that can be gained is dramatically different. The timing of returned signals can be converted into ranges, and readings are gathered across a spread of angles along a horizontal arc (azimuth). However, intensities are integrated over a vertical beam fan, meaning the elevation structure of the scene is collapsed into a flat image. Analogously to optical cameras, which provide width and height but not depth, a sonar image provides range and azimuth information but not altitudes.

\subsection{Imaging Sonar}
\label{sec:imaging_sonar_bg}
The acoustic pulse in FLS can be viewed as a discrete set of azimuthal beams, each composed of a continuous vertical fan of rays. \cref{fig:sonar_overview} visualizes this process, showing the polar representation of FLS. The sonar camera is considered to be at the origin, looking down the -X axis. The observed pixel intensities are determined by first–surface intersections along each beam, modulated by range–dependent transmission loss, noise, and multi-path and incidence effects. With a known speed of sound, the time-of-flight converts each return to range, reducing meter-scale ambiguity and providing pixel distances.

Sonar returns are typically dominated by diffuse and specular reflections, with relatively low ambient contributions. In natural settings, most surfaces can be well approximated using only diffuse reflections, while objects such as ship hulls and smooth rocks require specular consideration. In scenes with large, smooth, specular boundaries (\eg indoor tanks), specular reflections can produce multi-path returns, though this is less common in natural environments~\citep{Wang2023Multipath}. In practice, the signal is more likely to scatter, producing speckle. As such, compared to optics, sonar images typically have significant speckle noise (from scattering and interference) and low signal-to-noise ratio~\cite{Huang2020}.


\subsection{Sonar 3D Recovery}
Before modern learning-based pipelines, early work in sonar aimed to recover 3D shape directly from intensity using methods such as \textit{space carving} and \textit{shape from shading} (SfS). Space carving begins from a volumetric enclosure and iteratively removes voxels that fail multi-view consistency tests, leaving a volume that agrees with all views~\cite{Kutulakos1999SC}. Aykin and Negahdaripour~\cite{Aykin2015SC1, Aykin2017SC2} adapted this to sonar, but found it limited to simple shapes and heavily dependent on many viewpoints and rotations. SfS~\cite{Horn1970SfS} instead assumes a basic reflectance model, \eg Lambertian, and estimates surface normals or seafloor heights from brightness~\cite{Li1992SfS, Bikonis2013SfS1, Bikonis2013SfS2, Yakun2024SfS}. These approaches show that intensity encodes some 3D information, but only under restrictive assumptions and with extensive view coverage. Modern methods relax these assumptions by explicitly modeling the forward process or learning from data.

Sensor fusion, either using multiple sonar cameras or a combination of different sensor types, can also help resolve geometric ambiguities~\cite{Qadri2024AONeuS, Joe2019Fusion, Joe2022Fusion, Rho2025Fusion}, but it comes with increased hardware and calibration costs. In our work, we focus on 3D recovery from a single imaging sonar with a single view. DeBertoli \etal~\cite{DeBertoli2019ElevateNet} demonstrated that this is possible using a CNN trained to predict pixel-wise altitudes. Their method outperformed space carving, but required a large quantity of synthetic data, with results primarily demonstrated on simple objects.


\subsection{Differentiable Sonar Rendering}

Inverse rendering estimates scene parameters by matching rendered outputs to measurements. This optimization approach underlies differentiable rendering and NeRF-style methods for single- or multi-view reconstruction~\citep{Patow2003InvRendSurvey, Kato2020DiffRendSurvey, Mildenhall2020NeRF}. Recent work \citep{Qadri2022NeuSIS,Xie2022SSSDiff2, Bore2023SSSDiff, Reed2023SASDiff, Xie2022SSSDiff1,Xie2024BathymetryFLSDiff} has brought differentiable rendering to sonar. Qadri \etal~\citep{Qadri2022NeuSIS} demonstrated accurate 3D reconstruction of complex objects from multiple views using a volumetric rendering approach with neural signed distance functions (SDF), showing that physically-based forward models can achieve accurate shape recovery. However, their method depends on numerous views and is sensitive to complex scenes with multiple objects~\cite{Oliveira2025SES}. Additionally, representing the scene implicitly using neural networks increases the computational cost and fitting time. Comparatively, our method represents the seafloor as an explicit height field, allowing us to render a scene in under 100 milliseconds and fit a target image in tens of seconds. 

Forward models for elevation recovery have also been applied specifically for bathymetry, but are primarily limited to side-scan sonar (SSS) and synthetic aperture sonar (SAS)~\cite{Xie2022SSSDiff2, Bore2023SSSDiff, Reed2023SASDiff, Xie2022SSSDiff1}. Xie \etal~\cite{Xie2024BathymetryFLSDiff} introduced a self-supervised bathymetry framework that represents the seafloor as a neural heightmap (with multi-resolution hash encoding) and renders FLS intensities via a volumetric model that integrates along the elevation arc. Their method demonstrated accurate seafloor height recovery, but depends on many views along a surveying path. Unlike learned volumetric frameworks, we render images with an analytic, physics-based surface-scattering forward model.

\vspace{4pt}
\noindent\textbf{Our Approach.}
Existing methods for recovering 3D geometry from sonar typically rely on multiple views, target scenes with a single object, or require significant training data. Our focus is single–view FLS bathymetry using a physically grounded differentiable sonar simulator. Compared to other methods, we directly optimize a seafloor height field and fit observations using a differentiable physics approach, bypassing repeated neural-field evaluations during optimization while retaining adaptability across domains.